\subsection{Vector Fields}
\lecdate{(check) Lec 13 - Feb 24 (Week 7)}
we jump arnd a bit - doing chs 6/9 kinda together.

whats a vector field? morally, a smooth choice of tangent vec at every pt in $M$.
eg, for $M\subseteq\bR^{n}$, tangent vecs live in $T_{p}M\subseteq\bR^{n}$.

\begin{defn}
	a \textbf{vector field} is a smooth map $X:M\gto\bR^{n}$ s.t.
	$X(p)\in T_{p}M\subseteq\bR^{n}$.
\end{defn}
what does it mean on an abstract mfld? ``lets start by cheating".
v1: look at charts.

\begin{defn}
	a \textbf{vector field} $X$ on $M$ is a map $p\mto X(p)\in T_{p}M$, such that
	for each chart $\vphi$, we have $D_{p}\vphi\circ X:U\gto\bR^{n}$ is smooth.
\end{defn}
this uses the identification $T_{p}M\simeq\bR^{n}$ induced by chart.
but what is that thing between $D_{p}\vphi\circ X$?

\begin{defn}
	the \textbf{tangent bundle} $TM$ of $M^{m}$ is a smooth $2m$-mfld with:
	\begin{itemize}
		\item pt set $\bigsqcup_{p\in M}T_{p}M$
		\item (non-max'l) atlas w charts of the form:
			\begin{equation*}
				\left(\bigsqcup_{p\in U}T_{p}M,(\vphi\circ\pi,D\vphi)\right) \qquad
				(\vphi\circ\pi,D_{p}\vphi):\bigsqcup_{p\in U}T_{p}M\gto\bR^{m}\times\bR^{m}
			\end{equation*}
	\end{itemize}
	for charts $(U,\vphi)$ in $M$.
\end{defn}
check: $\pi$ is smooth. remarks:
\begin{enumerate}[(1)]
	\item the tangent bundle comes w a proj $\pi:TM\gto M$, given by $T_{p}M\mto p$.
	\item for each $p\in M$, the preimg (or fiber) $\pi^{-1}(p)$ is a vspace.
	\item \textit{local triviality:} for each $p\in M$, there's a nbhd $U\ni p$ s.t. the diagram commutes:
		[TODO: add diagram], where $(D\vphi,\pi)$ is a smooth map and
		linear on each $\pi^{-1}(q)$.
\end{enumerate}

\lecdate{Lec 15 - Mar 5 (Week 8)}
accidentally git yeeted tuesdays notes so those are gone forever

last time: a VF $X$ on $M$ defns a flow $\Phi:J\subseteq\bR\times M\gto M$ where $\dot{\Phi}_{t}(p)=X_{\Phi_{t}(p)}$,
for some $J\ni\set{0}\times M$ open.

\begin{defn}
	a VF is \textbf{complete} if $J=\bR\times M$.
\end{defn}
for instance, $x\p_{x}$ is complete w $\Phi_{t}(x)=e^{tx}$, but
$x^{2}\p_{x}$ is not (escapes to $\infty$ in finite time).

\begin{thm}
	if $X$ has cpt sppt, it is complete.
\end{thm} \

\begin{pf} \vspace{-0.2in}
	\begin{enumerate}[1)]
		\item if $p\notin\supp(X)$, $X_{p}=0$ implies $\Phi_{t}(p)=p$ for all $t$.
			if $p\in\supp(X)$, then $\Phi_{t}(p)\in\supp(X)$ when defnd.
		\item since $J$ open, $\fa\ep>0$, there is an open $(-\ep,\ep)\times U_{\ep}$, where:
			\begin{equation*}
				U_{\ep} \ = \ \set{p\in M:\Phi_{t}(p) \trm{ defnd for } t \in (-\ep,\ep)}
			\end{equation*}
			these cvr $M$. by cptness, take finite subcvr $\set{U_{\ep_{i}}}$.
			then $U_{\ep_{\min}}$=M.
		\item by flow property, $\Phi_{\ep_{\min}/2}$ defnd $\fa p\in M$.
			thus $\Phi_{n\ep_{\min}/2}(p)$ is defined on $t\geq0$,
			and $h<0$ can be handled the same way.
	\end{enumerate}
\end{pf}
summary: if $X$ a complete VF, its flow defns a 1-param. family of diffeos.

\begin{crll}
	any VF on a cpt mfld is complete.
\end{crll}
is the crll out of order? oh well. note:
\begin{enumerate}[i)]
	\item $\Phi_{t}\circ\Phi_{s}=\Phi_{t+s}$
	\item $\Phi(t,p)$ is a smooth fn on $\bR\times M$
\end{enumerate}
so $\Phi(t,p)$ a smooth fn on $\bR\times M$.

\begin{prop}
	if $\Phi_{t}$ satisfies i) and ii) above, its the flow of a VF $X$, where:
	\begin{equation*}
		X_{p}(f) \ = \ \frac{\d}{\d t}\bigg\rvert_{t=0}f(\Phi_{t}(p))
	\end{equation*}
\end{prop}
in other words, $X_{p}$ is the vel vec of $\gamma(t)=\Phi_{t}(p)$.

\begin{pf}
	first, $X$ satisfies the product rule:
	\begin{equation*}
		X_{p}(fg) \ = \ \frac{\d}{\d t}\bigg\rvert_{t=0}f(\Phi_{t}(p))g(\Phi_{t}(p))
	\end{equation*}
	its a VF. trust. it is also an integral curve of $X$, as:
	\begin{equation*}
		\dot{\Phi}_{t}(p) \ = \ \frac{\d}{\d s}\bigg\rvert_{s=0}\Phi_{s}(\Phi_{t}(p)) \ = \ X_{\Phi_{t}(p)}
	\end{equation*}
	sure.
\end{pf}

example: given $A\in M_{n}(\bR)$, let $\Phi_{t}:\bR^{m}\gto\bR^{m}$ be:
\begin{equation*}
	\Phi_{t}(x) \ = \ e^{tA}x \ = \ \left(\sum_{n=0}^{\infty}\frac{t^{n}}{n!}A^{n}\right)x
\end{equation*}
then the corresp. VF is:
\begin{equation*}
	X(f)(x) \ = \ \frac{\d}{\d t}\bigg\rvert_{t=0}f(e^{tA}x)
	\ = \ \nabla f\cdot\frac{\d}{\d t}\bigg\rvert_{t=0}e^{tA}x
	\ = \ \nabla f\cdot Ax
\end{equation*}
so in ``wishy-washy physics notation", we have $X=Ax\cdot\nabla$.

now, an alt defn of the lie bracket.
we start with, naturally, the pullback.

\begin{defn}
	given $F:M\gto N$, we denote by $F^{*}f=f\circ F$ the \textbf{pullback} of $f:N\gto\bR$.

	if $F$ is a diffeo, then for a VF $X$ on $N$, we can define $F^{*}X$ by:
	\begin{equation*}
		F^{*}X_{F^{-1}(p)}=TF^{-1}(X_{p})
	\end{equation*}
\end{defn} \

\begin{defn}
	the \textbf{lie derivative} of $f:M\gto\bR$ wrt a VF $X$ is given as:
	\begin{equation*}
		L_{X}f \ = \ \frac{\d}{\d t}\bigg\rvert_{t=0}\Phi_{t}^{*}f \ \in \ C^{\infty}(M)
	\end{equation*}
	the lie derivative of a VF $Y$ wrt $X$ is:
	\begin{equation*}
		L_{X}Y \ = \ \frac{\d}{\d t}\bigg\rvert_{t=0}\Phi_{t}^{*}Y
	\end{equation*}
\end{defn}
check: $L_{X}f=X(f)$.

\begin{thm}
	$L_{X}Y \ = \ [X,Y]$.
\end{thm}
